\documentclass[a4paper,10pt]{article}
\usepackage[left=1.5cm,top=2cm,right=2.5cm,bottom=4cm]{geometry}
\usepackage[utf8]{inputenc}
\usepackage[activeacute,spanish]{babel}
\usepackage{multirow}	% Para poder unir filas en las tablas
\usepackage{colortbl}	% Para colorear tablas
\usepackage{latexsym,amssymb,oldlfont,amsmath}
\usepackage{graphicx,graphics}
\usepackage{multicol}
\usepackage{color}
\usepackage{listings}
\usepackage{lastpage}
\usepackage{enumitem}
\usepackage{dsfont} %para poner numero reales


%\usepackage{fontspec}
%\setmainfont{Times New Roman}
\newcommand{\R}{{\mathbb R}}
\newcommand{\si}{{\mathcal S}^n}
\renewcommand{\theequation}{\thesection.\arabic{equation}} % Cambia formato salida numeracion equation
\makeatletter\@addtoreset {equation}{section} \makeatother % Tambien se usa
\renewcommand{\baselinestretch}{1.3}
\usepackage{amssymb}
\usepackage{hyperref}
\definecolor{midnightblue}{rgb}{0.1,0.1,0.44}


\newcommand{\dis}{\displaystyle}
\renewcommand{\theenumii}{\arabic{enumii}}

\def\Z{\mathbb{Z}}
\def\N{\mathbb{N}}
\def\R{\mathbb{R}}
\def\M{\mathcal{M}_n(\mathbb{R})}
\def\sin{\operatorname{sen}}
\def\e{\operatorname{e}}
\DeclareMathOperator{\arcsec}{arcsec}
\DeclareMathOperator{\arccot}{arccot}
\DeclareMathOperator{\arccsc}{arccsc}



\usepackage{etoolbox,fancyhdr,xcolor}
\fancyhf{} % sets both header and footer to nothing
\renewcommand{\headrulewidth}{0pt}
%\renewcommand{\footrulewidth}{0.5pt}
\pagestyle{fancy}
\headheight=60pt %para cambiar el tamaño del encabezadox
\fancyhead[L] %la "L" indica a la izquierda
{
\begin{minipage}{6cm}
\includegraphics[width=0.9\textwidth]{logo_udp.eps}
\end{minipage}
}

%\fancyhead[C]
%{
%\textbf{Cálculo II - Guía 1}
%}


%\fancyhead[R] %la "L" indica a la izquierda
%{

%}
%\begin{minipage}{4cm}
%\small %tamaño de letra pequeño
%{
%\textsf
%{
%Universidad de Chile\\
%...}}
%\end{minipage}}




\fancyfoot[C]{\thepage\ de \pageref{LastPage}}
%Página \thepage\ de \pageref{LastPage}}% Páginas Pares numeradas a la derecha
%\fancyfoot[LO]{\thepage} % Páginas Impares numeradas a la izquierda

\title{AyG-Guia1}

\begin{document}

\begin{center}
\emph{Álgebra y Geometría} \\
\emph{Guía Nº1 Ejercicios Resueltos} \\
\emph{Primer Semestre 2017} \\
\emph{Instituto de Ciencias Básicas}
\end{center}


\vspace{5mm}

\textbf{Ejercicio:} Demuestre la siguiente ley de absorción: $p \vee (p \wedge q) \equiv p $ \\

\textbf{Solución:} Se realiza una tabla de verdad.

\begin{center}
\begin{tabular}{|l|l|l|l|}
\hline
\multicolumn{1}{|c|}{$p$} & \multicolumn{1}{c|}{$q$} & \multicolumn{1}{c|}{$p \wedge q$} & \multicolumn{1}{c|}{$p\vee (p\wedge q)$} \\
\hline \hline
\multicolumn{1}{|c|}{V} & \multicolumn{1}{c|}{V} & \multicolumn{1}{c|}{V} & \multicolumn{1}{c|}{V} \\
\hline
\multicolumn{1}{|c|}{V} & \multicolumn{1}{c|}{F} & \multicolumn{1}{c|}{F} & \multicolumn{1}{c|}{V} \\
\hline
\multicolumn{1}{|c|}{F} & \multicolumn{1}{c|}{V} & \multicolumn{1}{c|}{F} & \multicolumn{1}{c|}{F} \\
\hline
\multicolumn{1}{|c|}{F} & \multicolumn{1}{c|}{F} & \multicolumn{1}{c|}{F} & \multicolumn{1}{c|}{F} \\
\hline
\end{tabular}
\end{center}

Observe que la columna asociada a $p \vee (p \wedge q)$ coincide con la columna asociada a $p$. Así, la proposición $p\vee (p \wedge q)$ es equivalente a la proposición $p$.
\vspace{5mm}







\textbf{Ejercicio:} Demuestre la siguiente proposición compuesta:
                    \[ \left[ \sim p \Rightarrow \sim (p \Rightarrow q)\right] \vee \left[ \left(p \wedge (p \Rightarrow q)\right) \Rightarrow p \right] \]
                    es una tautología

                    \begin{enumerate}
                        \item[a)] mediante álgebra de proposiciones.
                        \item[b)] mediante una tabla de verdad.
                    \end{enumerate}

\textbf{Solución:}

    \begin{enumerate}
        \item[a)] Simplificando:
            \begin{align*}
                       & \left[ \sim p \Rightarrow \sim (p \Rightarrow q)\right] \vee \left[ \left(p \wedge (p \Rightarrow q)\right) \Rightarrow p \right] \\
                \equiv & \left[  p \vee \sim (p \Rightarrow q)\right] \vee \left[ \left(p \wedge (\sim p \vee q)\right) \Rightarrow p \right] \\
                \equiv & \left[  p \vee \sim (\sim p \vee q)\right] \vee \left[ \left((p \wedge \sim p )  \vee (p \wedge q)\right) \Rightarrow p \right] \\
                \equiv & \underbrace{\left[  p \vee ( p \wedge \sim q)\right]}_{p \hspace{2mm}\text{(ley absorción)}} \vee \left[ \left( F  \vee (p \wedge q)\right) \Rightarrow p \right] \\
                \equiv & p \vee \left[ \left(p \wedge q\right) \Rightarrow p \right] \\
                \equiv &p \vee \left[ \sim \left(p \wedge q\right) \vee p \right] \\
                \equiv &p \vee \left[ \sim p \vee \sim q \vee  p \right] \\
                \equiv &p \vee \left[ (\sim p \vee  p) \vee \sim q  \right] \\
                \equiv &p \vee \left[ V \vee \sim q  \right] \\
                \equiv &p \vee V \\
                \equiv & V
            \end{align*}
            \newpage



        \item[b)] Tabla de verdad:

        \begin{center}
        \begin{tabular}{|l|l|l|l|l|l|l|l|l|}
            \hline
            \multicolumn{1}{|c|}{$p$} & \multicolumn{1}{c|}{$q$} & \multicolumn{1}{c|}{$\sim p$}  & \multicolumn{1}{c|}{$p \Rightarrow q$} & \multicolumn{1}{c|}{$\sim(p \Rightarrow q)$} & \multicolumn{1}{c|}{$\sim p \Rightarrow \sim (p \Rightarrow q)$ ($\ast$)} & \multicolumn{1}{c|}{$p \wedge (p\Rightarrow q)$} & \multicolumn{1}{c|}{$\left(p \wedge (p \Rightarrow q)\right) \Rightarrow p$ ($\ast \ast$)}  & \multicolumn{1}{c|}{($\ast$) $\vee$ ($\ast \ast$)} \\
            \hline
            \multicolumn{1}{|c|}{V} & \multicolumn{1}{c|}{V} & \multicolumn{1}{c|}{F} & \multicolumn{1}{c|}{V} & \multicolumn{1}{c|}{F} & \multicolumn{1}{c|}{V} & \multicolumn{1}{c|}{V} & \multicolumn{1}{c|}{V} & \multicolumn{1}{c|}{V} \\
            \hline
            \multicolumn{1}{|c|}{V} & \multicolumn{1}{c|}{F} & \multicolumn{1}{c|}{F} & \multicolumn{1}{c|}{F} & \multicolumn{1}{c|}{V} & \multicolumn{1}{c|}{V} & \multicolumn{1}{c|}{F} & \multicolumn{1}{c|}{V} & \multicolumn{1}{c|}{V} \\
            \hline
            \multicolumn{1}{|c|}{F} & \multicolumn{1}{c|}{V} & \multicolumn{1}{c|}{V} & \multicolumn{1}{c|}{V} & \multicolumn{1}{c|}{F} & \multicolumn{1}{c|}{F} & \multicolumn{1}{c|}{F} & \multicolumn{1}{c|}{V} & \multicolumn{1}{c|}{V} \\
            \hline
            \multicolumn{1}{|c|}{F} & \multicolumn{1}{c|}{F} & \multicolumn{1}{c|}{V} & \multicolumn{1}{c|}{V} & \multicolumn{1}{c|}{F} & \multicolumn{1}{c|}{F} & \multicolumn{1}{c|}{F} & \multicolumn{1}{c|}{V} & \multicolumn{1}{c|}{V} \\
            \hline
        \end{tabular}
        \end{center}

    \end{enumerate}
    \vspace{5mm}







\textbf{Ejercicio:} Simplifique las expresiones:
        \begin{enumerate}
            \item[a)] $\displaystyle \sim \left[ (p \Rightarrow q) \wedge (\sim p \Rightarrow q) \right]$
            \item[b)] $\displaystyle \left[ (p \Leftrightarrow \sim q) \wedge (q \wedge \sim p) \right]$
        \end{enumerate}

\textbf{Solución:}
        \begin{enumerate}
            \item[a)]
                \begin{align*}
                           & \sim \left[ (p \Rightarrow q) \wedge (\sim p \Rightarrow q) \right] \\
                    \equiv & \sim \left[ (\sim p \vee q) \wedge (p \vee q) \right] \\
                    \equiv & \sim \left[ (\sim p \wedge p) \vee q \right] \\
                    \equiv & \sim \left[ F \vee q \right] \\
                    \equiv & \sim q \\
                \end{align*}

            \item[b)]
                \begin{align*}
                           & \left[ (p \Leftrightarrow \sim q) \wedge (q \wedge \sim p) \right] \\
                    \equiv & \left[ \Big( (p \Rightarrow \sim q) \wedge (\sim q \Rightarrow p) \Big) \wedge q \wedge \sim p \right] \\
                    \equiv & (p \Rightarrow \sim q) \wedge (\sim q \Rightarrow p)  \wedge q \wedge \sim p \\
                    \equiv & (\sim p \vee \sim q) \wedge \underbrace{(q \vee p)  \wedge q}_{q} \wedge \sim p \\
                    \equiv & (\sim p \vee \sim q) \wedge q \wedge \sim p \\
                    \equiv & \underbrace{(\sim p \vee \sim q) \wedge \sim p}_{\sim p} \wedge q \\
                    \equiv & \sim p \wedge q
                \end{align*}
        \end{enumerate}
        \newpage







\textbf{Ejercicio:} Se sabe que la proposición $\displaystyle \left[ (q \Leftrightarrow p) \wedge \sim q \right] \Rightarrow (p\wedge \sim q) $ es falsa. Determine los valores de verdad de $p$ y $q$. \vspace{3mm}

\textbf{Solución:}

    Observe que la proposición consiste en una implicancia cuyo antecedente es $(q \Leftrightarrow p) \wedge \sim q$ y cuyo consecuente es $p \wedge \sim q$.

    El único caso en que la implicancia entrega un resultado falso es que el antecedente sea verdadero y el consecuente sea falso. Entonces,

    \begin{align*}
                        &  \left[ (q \Leftrightarrow p) \wedge \sim q \right] \Rightarrow (p\wedge \sim q) \equiv F \\
        \Leftrightarrow &  \left[ (q \Leftrightarrow p) \wedge \sim q \right] \equiv V  \hspace{2mm}  \wedge  \hspace{2mm}  (p \wedge \sim q) \equiv F \\
        \Leftrightarrow &  \left[ q \Leftrightarrow p \equiv V  \wedge \sim q \equiv V \right] \hspace{2mm}  \wedge  \hspace{2mm}  (p \wedge \sim q) \equiv F \\
        \Leftrightarrow &  \left[ q \Leftrightarrow p \equiv V  \wedge  q \equiv F \right] \hspace{2mm}  \wedge  \hspace{2mm}  (p \wedge \sim q) \equiv F \\
        \Leftrightarrow &  (q \equiv F) \wedge ( q \Leftrightarrow p \equiv V ) \wedge   (p \wedge \sim q \equiv F) \\
        \Leftrightarrow &  (q \equiv F) \wedge ( F \Leftrightarrow p \equiv V ) \wedge   (p \wedge \sim F \equiv F) \\
        \Leftrightarrow &  (q \equiv F) \wedge ( p \equiv F ) \wedge   (p \wedge V \equiv F) \\
        \Leftrightarrow &  (q \equiv F) \wedge ( p \equiv F )
    \end{align*}

    Luego, tanto $p$ como $q$ son falsas.
    \vspace{5mm}








\textbf{Ejercicio:} Se sabe que $\left[(p \wedge \sim q) \Rightarrow r \right] \Rightarrow (\sim p \vee \sim s)$ es falsa, y que $(\sim s \vee \sim r)$ es verdadera. Encuentre los valores de verdad de $p,q,r$ y $s$. \vspace{3mm}

\textbf{Solución:}

    De la primera proposición, se deduce que $(p \wedge \sim q) \Rightarrow r \equiv V$ y que $(\sim p \vee \sim s) \equiv F$.

    Si $(\sim p \vee \sim s) \equiv F$ entonces $p \wedge s \equiv V$, por lo que $p\equiv V$ y $s\equiv V$.

    Por otro lado, considerando que $s \equiv V$, la segunda proposición entrega como información que $(\sim V \vee \sim r) \equiv V$, es decir $\sim r \equiv V$, por lo que $r \equiv F$.

    Finalmente, considerando estos valores de verdad, en conjunto con la primera proposición, se tiene que $(V \wedge \sim q) \Rightarrow F \equiv V$, es decir $\sim q \Rightarrow F \equiv V$. En otras palabras, $\sim q \equiv F$, y así $q \equiv V$.

    En conclusión: $p,q,s$ son verdaderas, y $r$ es falsa.

\end{document}
